{%- macro bold_red(character) -%}
    \textcolor{red}{\textbf{ {{character}} }}
{%- endmacro -%}

{% macro pop_set(set_) %}
    {%- if not set_ -%}
    \varnothing
    {%- else -%}
        \{ 
        {%- for item in set_ -%}
            {{ item }}
           {%- if not loop.last -%}
            , 
           {%- endif -%}
        {%- endfor -%}
        \}
    {%- endif %}
{% endmacro %}

\section{Removing $\epsilon$-Transitions}

Our goal in this section is to eliminate all $\epsilon$-transitions from the given NFA while preserving the language it accepts. Intuitively, an $\epsilon$-transition allows the automaton to change states without consuming any input symbol. While this is convenient for construction, it complicates execution and analysis. To address this, we construct a new automaton $M'$ that has no $\epsilon$-transitions but behaves \emph{exactly the same} as the original automaton $M$ on every input string.

\textbf{Key idea:} Any time the original NFA could take one or more $\epsilon$-transitions before or after reading an input character, we explicitly encode that behavior using ordinary transitions in $M'$.

\begin{figure}[H]
    \centering
    \begin{minipage}{0.47\linewidth}
        \centering
        \includegraphics[width=0.95\linewidth]{figures/original.png}
    \end{minipage}\hfill
    \begin{minipage}{0.47\linewidth}
        \centering
        \includegraphics[width=0.95\linewidth]{figures/final_no_epsilon.png}
    \end{minipage}
    \caption{Original NFA (Left) And NFA Without $\epsilon$-Transitions (Right)}
\end{figure}

The automaton on the right has no $\epsilon$-transitions, but every path that was possible in the original automaton is still represented.

\subsection*{$\epsilon$-Closure Intuition}

When starting from a state in an NFA, it is often possible to reach several other states without reading any input symbols by following $\epsilon$-transitions. This means that the automaton may \emph{effectively} be in multiple states at once before reading the next character. To capture this behavior, we use the concept of the \textbf{$\epsilon$-closure}. The $\epsilon$-closure of a set of states consists of all states that can be reached from that set by following zero or more $\epsilon$-transitions.

\textbf{Important observations:}
\begin{itemize}[itemsep=0.3em]
    \item The $\epsilon$-closure is always defined on a \emph{set} of states, not just a single state.
    \item Every state is always in its own $\epsilon$-closure.
    \item $\epsilon$-closures capture all ``hidden movement'' the automaton can make without reading input.
\end{itemize}

{%- if e_closure_table %}
\begin{center}
\setlength{\tabcolsep}{12pt}
\renewcommand{\arraystretch}{1.1}
\begin{tabular}{c|l}
\textbf{State} & $\epsilon$-closure \\\hline
{%- for state, closure in e_closure_table.items() %}
{{ state }} & $\{ {{ closure }} \}$ \\\hline
{%- endfor %}
\end{tabular}
\end{center}
{%- endif %}

The table above shows the $\epsilon$-closure for each state. For example, if $E(\{q\}) = \{q, r, s\}$, then starting in $q$ the automaton could silently move to $r$ or $s$ before consuming any input.

---

{% for state, closure in e_closure_table.items() %}
\subsection*{Handling State ${{ state }}$}

We now describe how transitions in $M'$ are constructed for state ${{ state }}$. Each step ensures that
all possible behaviors involving $\epsilon$-transitions in $M$ are faithfully simulated.

\begin{enumerate}[itemsep=0.4em]

    \item \textcolor{gray}{Calculate the $\epsilon$-closure of the state in $M$; call the result $S$.}  

    From the table above, we know that:
    \[
        E(\{ {{state}} \}) = \{ {{closure}} \}.
    \]

    This means that before reading any input symbol, the automaton could already be in \emph{any} of the
    states in $\{ {{closure}} \}$.

    \item \textcolor{gray}{Calculate transitions from all states in $S$ on each input symbol.}

    For each input symbol, we collect every possible destination reachable by reading that symbol from
    \emph{any} state in the $\epsilon$-closure. This reflects the idea that the automaton might have silently
    moved before reading the symbol.

    \begin{itemize}[itemsep=0.2em]
        {% for character, value in transitions.get(state, 'ERROR').items() %}
            \item [{{bold_red(character)}}] 
            Starting in any state in $\{ {{closure}} \}$, we consider all transitions labeled
            {{bold_red(character)}}.

            \begin{itemize}
            {%- for e_closure_state in closure -%}
                \item $\delta({{e_closure_state}}, {{bold_red(character)}}) = {{ pop_set(transitions.get(e_closure_state, empty_set).get(character, empty_set)) }}$
            {%- endfor %}
            \end{itemize}

            Collecting these results, we obtain an intermediate set of states reachable after reading
            {{bold_red(character)}}:
            \[
                \delta_R({{state}}, {{bold_red(character)}}) \supseteq
                \{ {{partial_closure.get(state, empty_set).get(character, empty_set)}} \}.
            \]
        {% endfor %}
    \end{itemize}

    \item \textcolor{gray}{Take the $\epsilon$-closure of the resulting set.}

    After reading an input symbol, the automaton may once again follow $\epsilon$-transitions before
    processing the next symbol. Therefore, we compute the $\epsilon$-closure of the intermediate set to
    obtain the final transition in $M'$.

    \begin{itemize}[itemsep=0.2em]
    {% for character, value in transitions.get(state, 'ERROR').items() %}
        \item [{{bold_red(character)}}]
        Taking the $\epsilon$-closure of the intermediate set, we conclude:
        \[
            \delta_R({{state}}, {{character}}) =
            \{ {{aggregate_closure.get(state, empty_set)[character]}} \}.
        \]
    {% endfor %}
    \end{itemize}

\end{enumerate}

After processing state {{state}}, we add these transitions to $M'$. The updated automaton is shown below.

\begin{figure}[H]
    \centering
    \vspace{0.3em}
    \includegraphics[width=0.82\linewidth]{figures/step_{{state}}.png}
    \vspace{0.3em}
    \caption{Updated Automaton After Handling State ${{ state }}$}
    \label{fig:M_after_step_{{state}}}
\end{figure}

{% endfor %}

\subsection*{Final Result}

By repeating this process for every state in the original automaton, we ensure that all behaviors involving
$\epsilon$-transitions are fully simulated using only ordinary transitions.

\textbf{Conclusion:} The resulting automaton $M'$ has no $\epsilon$-transitions and accepts exactly the
same language as the original NFA $M$. All implicit movement has been made explicit, simplifying
execution while preserving correctness. 
